\section{Synthesis methodology: five routes to the same framework}
\label{sec:synthesis}

The final lever (L5) is often treated as logistics, but the library evidence says otherwise:
the synthesis route sets the crystallinity and phase purity on which every upstream factor
depends~\cite{ghosh2020identification,zhang2022reconstructed}. This section compares the five
route families documented for the TpPa system, the target chemistry of Route~C.

\textbf{Solvothermal baseline and solvent thermodynamics (L5a).} The canonical TpPa-1 protocol
--- Tp + Pa-1 in 1:1 mesitylene/dioxane with aqueous acetic acid, sealed at 120~$^\circ$C for
72~h --- remains the reference for crystallinity and chemical
robustness~\cite{kandambeth2012construction}. Solvent choice, long empirical, now has a
quantitative decision framework: ranking solvents by CHEM21 greenness, Hansen solubility
parameters, and free energy of solvation shows that moderate-solvation media (1,4-dioxane,
propylene carbonate, diacetin) suppress excessive reaction reversibility and promote
crystallite growth~\cite{xu2026structural}. New synthetic strategies across the COF field are
compiled in the dedicated CSR review~\cite{li2020new}.

\textbf{Rapid and green routes (L5b).} Four alternatives trade conditions against quality:
\begin{itemize}
\item \emph{Mechanochemical grinding}: room-temperature, solvent-free, delivering the same
  chemical stability with moderate crystallinity and partially exfoliated
  morphology~\cite{biswal2013mechanochemical}.
\item \emph{Salt-mediated (``organic terracotta'')}: p-toluenesulfonic acid mediation
  crystallizes ultraporous frameworks in seconds, with family surface areas up to
  3000~m$^2$\,g$^{-1}$ and continuous processing demonstrated on a twin-screw
  extruder~\cite{karak2017constructing}.
\item \emph{Microwave-assisted solvothermal}: an enamine COF forms in significantly less time
  at lower temperature with high yield~\cite{wei2015the}; in the direct TpPa-1 comparison
  compiled by the microwave review, 1~h at 100~$^\circ$C gives a BET surface area of
  725~m$^2$\,g$^{-1}$ versus 152~m$^2$\,g$^{-1}$ for the conventional 3-day, 120~$^\circ$C
  control~\cite{grenu2020microwave}. The same review documents an instructive caveat: for
  transimination-based crystallinity upgrading, solvothermal heating still outperforms
  microwave for $\beta$-ketoenamine COFs --- the ``fastest route'' is
  context-dependent~\cite{grenu2020microwave}.
\item \emph{Room-temperature and flow synthesis}: RT batch TpPa-1 matches solvothermal
  crystallinity and porosity, attributed to monomer solubility and $\pi$-stacking-assisted
  error correction; the first continuous-flow COF production reached 41~mg\,h$^{-1}$ at a
  space-time yield of 703~kg\,m$^{-3}$\,day$^{-1}$~\cite{peng2016room}. Flow synthesis of
  TpPa-1 in the green solvent diacetin reaches a BET area of 418~m$^2$\,g$^{-1}$ with a
  30-fold space-time-yield increase and an 89\% reduction in specific energy relative to batch
  in the same solvent~\cite{xu2026structural}.
\end{itemize}

\textbf{High-throughput discovery (L5c).} Scale-out of synthesis itself is emerging:
sonochemical high-throughput experimentation screened 76 polymers, of which 60 formed
crystalline COFs and 18 were previously unreported, in a single campaign --- the target
application was photocatalytic H$_2$O$_2$ production, but the accelerated
synthesis-plus-screening methodology transfers directly to HER~\cite{zhao2022accelerated}.
The campaign format matters for this survey's argument: it is the first library-scale
demonstration that route parameters (here, sonication) can be co-screened with photocatalytic
output, rather than fixed by convention before testing begins~\cite{zhao2022accelerated}.
Combined with AQE-standardized reporting~\cite{li2022covalent} and factor
benchmarking~\cite{ghosh2020identification}, this points toward synthesis--performance maps
built at campaign scale rather than paper scale.

\begin{table}[t]
\centering
\caption{Synthesis routes to TpPa-family COFs (taxonomy level L5). Quantities are as reported
in the cited sources; dashes mark values not reported in the library evidence.}
\label{tab:routes}
\footnotesize
\begin{tabular}{p{2.7cm}p{3.1cm}p{3.5cm}p{3.6cm}}
\toprule
Route & Conditions & Reported quality & Scale/greenness notes \\
\midrule
Solvothermal~\cite{kandambeth2012construction} & 120~$^\circ$C, 72~h,
mesitylene/dioxane + aq.~AcOH & Reference crystallinity; 9~N HCl and boiling-water stability &
Sealed-tube batch; high-boiling solvents \\
Mechanochemical~\cite{biswal2013mechanochemical} & Room temperature, solvent-free grinding &
Moderate crystallinity; stability retained; exfoliated layers & Greenest inputs; lowest
crystallinity \\
Salt-mediated~\cite{karak2017constructing} & p-TsOH mediation, seconds & Highly crystalline,
ultraporous (up to 3000~m$^2$\,g$^{-1}$, family value) & Twin-screw extrusion; salt work-up
required \\
Microwave~\cite{wei2015the,grenu2020microwave} & 100~$^\circ$C, 1~h & BET 725 vs.\
152~m$^2$\,g$^{-1}$ (conventional control) & Fastest lab route; penetration-depth scale limit \\
RT / flow~\cite{peng2016room,xu2026structural} & RT batch; flow in diacetin & RT: comparable to
solvothermal; flow: BET 418~m$^2$\,g$^{-1}$ & Flow: 30$\times$ STY, $-$89\% specific energy;
COF flow precedent 703~kg\,m$^{-3}$\,day$^{-1}$ \\
\bottomrule
\end{tabular}
\end{table}

The pattern across Table~\ref{tab:routes} motivates the design task of the next section: no
single route dominates all axes, and the HER consequences of the green/scalable routes are
essentially unmeasured (Sec.~\ref{sec:challenges}, G4). Choosing a route is therefore a
decision problem with documented trade-offs --- exactly the kind of problem a survey should
hand over to synthesis planning in executable form.
